\documentclass[11pt,a4paper]{article}
\usepackage[margin=1in]{geometry}
\usepackage{hyperref}
\usepackage{enumitem}
\usepackage{graphicx}
\usepackage{xcolor}

% -----------------------------------------------------
% bvraghav-tiet-ucs503p-article.sty
% -----------------------------------------------------

% Variable: Lab Instructor
\makeatletter \newcommand{\BvrSetLabInstructor}[1]{%
  \gdef\Bvr@LabInstructor{#1}}
% default
\newcommand{\Bvr@LabInstructor}{Ms. Paramveer Sidhu}
% public accessor
\newcommand{\BvrLabInstructorValue}{\Bvr@LabInstructor}
\makeatother

% Add subtitle
\usepackage{titling}
% -- Set title bold
\pretitle{\begin{center}\bfseries\fontsize{18bp}{18bp}\selectfont}
\makeatletter
\newcommand{\BvrSubtitle}[1]{%
  \posttitle{%
    \par\end{center}
    \begin{center}\large#1\end{center}
    \vskip 0.5em}%
}
\makeatother

% Hints in the section title(s)
\newcommand{\BvrHint}[1]{%
  {\normalfont\normalsize\itshape\hfill #1}}

% -----------------------------------------------------

\title{UniVerse: A Trust-Based Campus Community and Skill-Exchange Platform}

\BvrSetLabInstructor{Ms. Paramveer Sidhu}

\BvrSubtitle{%
  Project Proposal for UCS503P  \\
  Submitted to: \BvrLabInstructorValue \\ \bfseries
  Thapar Institute of Engineering and Technology% 
}

\author{
  Sehajneet Kaur, CSED%
  \thanks{Roll No: \texttt{1024030215}, %
    Email: \texttt{skaur10\_be24\@thapar.edu}}
  \and 
  Raynee Jindal, CSED%
  \thanks{Roll No: \texttt{1024030711}, %
    Email: \texttt{rjinadal\_be24\@thapar.edu}}
  \and
  Sahejbir Singh, CSED%
  \thanks{Roll No: \texttt{1024030247}, %
    Email: \texttt{ssingh10\_be2x\@thapar.edu}}
  \and
  Shaurya Ranjan, CSED%
  \thanks{Roll No: \texttt{1024030694}, %
    Email: \texttt{sranjan\_be24\@thapar.edu}}
}

\date{\today}

\begin{document}
\maketitle

\section{Higher-order goal%
  \BvrHint{\ldots secondary framing}}
This project aims to strengthen peer-to-peer trust and knowledge circulation within a college campus by making mentorship, skill-sharing, and resource-lending discoverable to every student, not just those already embedded in a society or a large personal network. While the broader goal is community-building, the immediate engineering objective is to translate \emph{trust-based peer exchange} into a measurable, deployable software solution.

\section{Time-to-value %
\BvrHint{\ldots primary heuristic}}
\begin{itemize}[leftmargin=0.7cm]
    \item We will deliver a usable core loop quickly by prototyping the channel-based community space and the skill-swap matching flow, and validating both with real students within a short iteration window.
    \item We will prioritize fast feedback over perfection by using weekly iterations and targeting CI-backed automation early to reduce release friction.
    \item Success is defined by what can be delivered and measured in the first iteration -- a working set of community channels and a functioning skill-swap points system -- then improved based on user feedback.
\end{itemize}

\section{Problem Statement}
Students at a residential college campus often cannot access mentorship, specific skills, or basic resources (borrowing an item, recovering something lost) unless they already belong to a society or have a wide personal network. Common pain points include:
\begin{itemize}[leftmargin=0.7cm]
    \item \textbf{Access inequality:} guidance and mentorship on niche or specific skills is concentrated within existing social circles, leaving unconnected students without a path in.
    \item \textbf{Redundant paid learning:} students often pay for external resources to learn something a peer on campus could already teach them.
    \item \textbf{No trust layer for informal exchange:} lending, borrowing, lost-and-found, and general help currently happen through scattered WhatsApp groups with no accountability, history, or verification.
    \item \textbf{No record of informal learning:} peer-taught skills and mentorship sessions leave no trace a student can point to later, even though the effort and outcome are real.
\end{itemize}

\textbf{Why this matters now (contextual relevance):} With a large, skill-diverse student body living on one campus, the supply of mentorship and resources already exists -- what is missing is a trusted, discoverable layer to connect it.

\section{Proposed Solution %
\BvrHint{\ldots Software Proposition}}
\subsection{Overview}
Build a \textbf{web-based} campus community platform, UniVerse, with two core pillars:
\begin{itemize}[leftmargin=0.7cm]
    \item \textbf{Community channels:} Discord-style dedicated channels for lending, borrowing, lost items, found items, and general discussion.
    \item \textbf{Skill Swap:} a mentorship and peer-learning marketplace where users list skills they can teach and skills they want to learn, request sessions, and build a points-based reputation.
    \item \textbf{User profile:} name, branch, skills offered/wanted, points balance, sessions taken/given, and leaderboard rank -- shareable as proof of learning (e.g., on LinkedIn).
    \item \textbf{Recommendation surface:} the Skill Swap landing page recommends listings based on the skills a user has marked as wanted.
\end{itemize}

\subsection{Core workflow}
\begin{enumerate}[leftmargin=0.7cm]
    \item Student creates a profile listing skills they can teach and skills they want to learn.
    \item Skill Swap surfaces recommended mentors/listings based on the student's "want to learn" list.
    \item Student sends a request for a session (guidance, mentorship, or skill swap); once completed, the mentor earns points and the learner spends points.
    \item Points scale with uniqueness of mentees taught, feeding a leaderboard that doubles as a portable record of contribution.
    \item In parallel, students post and browse lending, borrowing, lost, and found items in dedicated channels, building a visible history of trustworthy exchanges.
\end{enumerate}

\subsection{Operational constraints for an educational, campus setting}
\begin{itemize}[leftmargin=0.7cm]
    \item Keep the solution usable on low-to-mid range devices via responsive web design.
    \item Avoid dependence on advanced research algorithms (e.g., ML-based matching); focus on workflow, trust mechanics, and system correctness.
    \item Restrict sign-up to verified college email addresses to keep the trust layer meaningful.
    \item Design for deployability using common web hosting and managed databases.
\end{itemize}

\section{Solution Approach %
\BvrHint{\ldots Engineering focus}}
This is an engineering course project, so we focus on scalable platform and workflow aspects rather than novel algorithm research.

\subsection{Points and reputation logic %
\BvrHint{\ldots non-research}}
The points system will be rule-based, not ML-based:
\begin{itemize}[leftmargin=0.7cm]
    \item Teaching a skill awards more points than learning one costs, by a fixed configurable ratio.
    \item Teaching additional \emph{unique} learners yields a bonus, discouraging repeated sessions with the same pair from inflating rank.
    \item All point transactions are logged per session for auditability and dispute resolution; there is no automatic deduction without a completed, confirmed session.
\end{itemize}

\subsection{Recommendation logic}
Listing recommendations on the Skill Swap page will use a simple tag-matching heuristic (overlap between a user's "want to learn" tags and other users' "can teach" tags), without claiming state-of-the-art personalization.

\subsection{Web architecture %
\BvrHint{\ldots web-based solution}}
A typical 3-tier web architecture:
\begin{itemize}[leftmargin=0.7cm]
    \item \textbf{Frontend:} React-based responsive UI for channels, profiles, Skill Swap listings, and the leaderboard.
    \item \textbf{Backend API:} Node.js/Express service handling authentication, channel messaging, session requests, and points transactions.
    \item \textbf{Database:} MongoDB storing users, profiles, channel messages, skill listings, sessions, and points history.
\end{itemize}

\subsection{CI/CD and fast delivery enabler}
CI/CD will be used to:
\begin{itemize}[leftmargin=0.7cm]
    \item Automatically build and run tests on every change.
    \item Produce repeatable deployment artifacts to staging.
    \item Enable safe and frequent releases of incremental features (skill swap first, then channels, then leaderboard/profile polish).
\end{itemize}

\section{Evaluation Criterion %
\BvrHint{\ldots measurable and attributable}}
We define success using measurable metrics tied to the two core workflows.

\subsection{Primary evaluation metric}
\textbf{Session completion rate:} proportion of Skill Swap session requests that result in a confirmed, completed session.
\begin{itemize}[leftmargin=0.7cm]
    \item Target: $\ge$ 60\% of session requests completed during the pilot window.
    \item Attribution: measured from database timestamps and session-confirmation records (request created vs. marked complete by both parties).
\end{itemize}

\subsection{Secondary evaluation metrics}
\begin{itemize}[leftmargin=0.7cm]
    \item \textbf{Unique mentee reach:} average number of distinct learners taught per active mentor.
    \item \textbf{Channel activity:} number of lending/borrowing/lost/found posts resolved (item returned or matched) versus posted.
    \item \textbf{User clarity:} user-reported clarity of the points and ranking system using a simple 1--5 feedback question.
    \item \textbf{Reliability:} availability of staging/production for login, posting, and session requests during peak campus hours (e.g., $\ge$ 99\% uptime in pilot).
\end{itemize}

\subsection{Pilot validation plan %
\BvrHint{\ldots high-level}}
\begin{itemize}[leftmargin=0.7cm]
    \item Recruit 20--30 pilot users across multiple branches and years to ensure skill diversity.
    \item Track requested vs. completed sessions and channel resolutions over the iteration window, comparing against a baseline collected via short interviews on how students currently find mentors or lend/borrow items.
\end{itemize}

\section{Scalability %
\BvrHint{\ldots with foresight}}
The proposition should scale in both theoretical and practical senses.

\begin{enumerate}[leftmargin=0.7cm]
    \item \textbf{Scalable (theoretically):} The system should support growth in number of users, channels, and concurrent Skill Swap requests by:
    \begin{itemize}[leftmargin=0.7cm]
        \item decoupling components (frontend/backend),
        \item enabling pagination and indexing for channel messages and skill listings,
        \item using stateless API design.
    \end{itemize}
    
    \item \textbf{Operationally deployable (practically):} The system should run on common web hosting:
    \begin{itemize}[leftmargin=0.7cm]
        \item managed MongoDB instance,
        \item containerized or platform-hosted deployment,
        \item CI/CD pipeline for staging/production.
    \end{itemize}
\end{enumerate}

\section{Engine availability heuristic}
A mature set of ``engines'' (tooling/services) is available to implement this as a web-based project:
\begin{itemize}[leftmargin=0.7cm]
    \item React for the frontend; Node.js/Express for REST APIs.
    \item MongoDB (managed, e.g., Atlas) as the primary document store.
    \item Token-based authentication restricted to college email domains.
    \item Object storage for optional attachments (item photos) or local storage for pilot.
    \item Test framework for unit/integration tests (e.g., Jest).
    \item CI runner and staging deployment target.
\end{itemize}
We will use these mature components to focus on workflow design, trust mechanics, and measurement rather than inventing new infrastructure.

\section{Project scope and deliverables %
\BvrHint{\ldots iteration-friendly}}
\subsection{Initial deliverable %
\BvrHint{\ldots first iteration}}
\begin{itemize}[leftmargin=0.7cm]
    \item User authentication (college email) and basic profile (name, branch, skills)
    \item Community channels: lending, borrowing, lost, found, general
    \item Skill Swap: list skills to teach/learn, browse listings
    \item Session request flow with manual completion confirmation
    \item Points logic: award/deduct on confirmed session completion
    \item CI: build + backend unit tests + linting
    \item CD to staging with a single command or pipeline job
\end{itemize}

\subsection{Subsequent deliverables}
\begin{itemize}[leftmargin=0.7cm]
    \item Tag-based recommended listings on the Skill Swap landing page
    \item Leaderboard with rank, sessions taught/taken, and unique-mentee bonus
    \item Shareable profile summary (exportable for LinkedIn)
    \item Improved search/filtering and indexed queries for scale readiness
\end{itemize}

\section{Risks and mitigations}
\begin{itemize}[leftmargin=0.7cm]
    \item \textbf{Trust and safety:} restrict sign-up to verified college emails; keep session and item-exchange history visible on profiles to discourage bad behavior.
    \item \textbf{Points gaming:} require both-party confirmation before awarding/deducting points; cap bonus effects; log all transactions for audit.
    \item \textbf{Low initial supply of mentors:} seed the platform with a founding cohort of students across branches before wider launch.
    \item \textbf{Evaluation measurement ambiguity:} instrument timestamps and define clear session/channel-resolution states before development begins.
    \item \textbf{Operational issues in pilot:} use staging early; ensure CI/CD produces repeatable deployments.
\end{itemize}

\section{Summary}
This project proposes UniVerse, a deployable campus web platform that reduces the friction of finding mentorship, learning niche skills, and trusting peer-to-peer lending and lost-and-found exchange. It meets the course criteria by emphasizing time-to-value through rapid iterations, implementing CI/CD to support frequent safe releases, and using measurable evaluation criteria (especially session completion rate). It remains focused on engineering workflow, trust mechanics, and scalability readiness rather than research-driven novelty.

\end{document}
